% Clase 4 — conductor en equilibrio con hueco sin cargas.
% (a) gaussiana S pegada al hueco, del lado del material: E=0 sobre S => Q_int=0.
% (b) por qué la carga neta nula no basta: parches +/- compensados y la curva C
%     cerrada que los descarta (circulacion no nula).
\begin{tikzpicture}[scale=1]

% ---------------- panel (a) ----------------
\begin{scope}
  \fill[figgray!18] (0,0) circle (2.05);
  \draw[figline] (0,0) circle (2.05);
  \fill[white] plot[smooth cycle, tension=0.9]
    coordinates {(-0.75,0.55) (0.15,0.9) (0.9,0.3) (0.6,-0.65) (-0.4,-0.6)};
  \draw[figline] plot[smooth cycle, tension=0.9]
    coordinates {(-0.75,0.55) (0.15,0.9) (0.9,0.3) (0.6,-0.65) (-0.4,-0.6)};
  % superficie gaussiana S, pegada al hueco pero del lado del material
  \draw[figred, dashed, line width=0.9pt] plot[smooth cycle, tension=0.9]
    coordinates {(-1.0,0.75) (0.2,1.18) (1.15,0.4) (0.82,-0.9) (-0.6,-0.85)};
  \node[figlbl, figred, fill=white, inner sep=1.5pt] at (1.02,1.28) {$S$};
  \node[figlbl] at (0.08,0.1) {hueco};
  \node[figlblaux, fill=figgray!18, inner sep=1.5pt] at (0,-1.5) {$\vec E=0,\ \rho=0$};
  \node[figlbl] at (0,-2.55) {conductor en equilibrio};
  \node[figlbl] at (0,-3.15) {(a) $\vec E=0$ sobre $S\ \Rightarrow\ Q_{\text{sup. int}}=0$};
\end{scope}

% ---------------- panel (b) ----------------
\begin{scope}[xshift=7.4cm]
  \fill[figgray!18] (0,0) circle (2.05);
  \draw[figline] (0,0) circle (2.05);
  \fill[white] plot[smooth cycle, tension=0.85]
    coordinates {(-1.3,0.5) (-0.1,0.85) (1.2,0.5) (1.25,-0.2) (-0.1,-0.55) (-1.3,-0.2)};
  \draw[figline] plot[smooth cycle, tension=0.85]
    coordinates {(-1.3,0.5) (-0.1,0.85) (1.2,0.5) (1.25,-0.2) (-0.1,-0.55) (-1.3,-0.2)};
  % parches de carga sobre la pared del hueco
  \foreach \yy in {-0.1,0.06,0.22}{
    \node[figpos] at (-1.16,\yy) {};
    \node[figneg] at (1.11,\yy) {};
  }
  \node[figlbl, figred, fill=figgray!18, inner sep=1.5pt] at (-1.38,1.08) {$\sigma>0$};
  \node[figlbl, figblue, fill=figgray!18, inner sep=1.5pt] at (1.38,1.08) {$\sigma<0$};
  \draw[figvecaux] (-1.3,0.88) -- (-1.18,0.4);
  \draw[figvecaux] (1.3,0.88) -- (1.18,0.4);
  % campo dentro del hueco: de + a -
  \draw[figvec] (-0.88,0.08) -- (0.78,0.08);
  \node[figlbl] at (-0.05,0.36) {$\vec E$};
  % curva cerrada C: tramo por el hueco y retorno por el material
  \draw[figred, line width=1pt, dashed]
    (-0.95,-0.08) -- (0.82,-0.08) -- (0.82,-1.5) -- (-0.95,-1.5) -- cycle;
  \node[figlbl, figred, fill=figgray!18, inner sep=1.5pt] at (-0.05,-1.5) {$C$};
  \node[figlblaux, fill=figgray!18, inner sep=1.5pt] at (-0.05,-0.9) {aqu\'i $\vec E=0$};
  \node[figlbl] at (0,-2.55) {carga neta nula, pero $\sigma\neq0$};
  \node[figlbl] at (0,-3.15) {(b) $\oint_C\vec E\cdot d\vec\ell\neq0$: imposible};
\end{scope}

\end{tikzpicture}
